% SHARD-ID: AQM-64-QUANTUM-SYSTEM-ASSOCIATION-META-PLAN
% SHARD-TITLE: Quantum-System Association Meta-Plan
% SHARD-SUMMARY: Records a user-directed programme for cataloguing ways to associate quantum kinematical systems to objects.
% SHARD-SUMMARY: Reserves a sequence of review and new shards for finite sets, geometric objects, finite rings, agreement tests, and degeneracies.
% SHARD-SUMMARY: Points to the detailed plan and Beads chain without asserting new mathematical comparison results.
% SHARD-KEYWORDS: meta-plan, association, finite sets, finite rings, Weyl-Heisenberg, cotangent, review shards

\section{Quantum-System Association Meta-Plan}
\label{sec:quantum-system-association-meta-plan}

\subsection{Status}

This is a planning shard.  It records a user-directed investigation programme,
not a mathematical result.  Its detailed working document is
\texttt{docs/quantum\_system\_associations/PLAN.md}.  The persistent work
items are the Beads chain \texttt{aqm-qsa01} through \texttt{aqm-qsa23}.

The programme asks for a catalogue of ways to associate quantum kinematical
systems to an object \(X\).  The guiding convention is to avoid treating this
as a canonical functorial quantisation claim.  Each association workflow must
record the structure placed on \(X\), the extra choices made, the quantum
carrier or observable algebra produced, and the evidence supporting any
comparison with another workflow.

\subsection{Reserved Shard Sequence}

Each item below is intended to become one report shard.  A review shard means
that earlier report material already contains much of the evidence, but the new
shard must still cite the exact local anchors and explain the result in the
language of this meta-plan.

\begin{center}
\begin{tabular}{rlll}
\toprule
Step & Planned shard & Kind & Topic \\
\midrule
1 & AQM-65 & New & Association vocabulary and non-functoriality \\
2 & AQM-66 & New & Test object catalogue \\
3 & AQM-67 & New & Structureless finite set: first association \\
4 & AQM-68 & Review/New & Structureless finite set: tensor sites \\
5 & AQM-69 & New & Structureless finite set: direct-sum particle \\
6 & AQM-70 & New & AND/OR many-particle grammar \\
7 & AQM-71 & Review & Finite symplectic field labels on a set \\
8 & AQM-72 & Review/New & Bosonic and fermionic association layer \\
9 & AQM-73 & Gap & Fusion-category endpoint \\
10 & AQM-74 & New & Single-particle reduction \\
11 & AQM-75 & Review/New & What counts as a point \\
12 & AQM-76 & New & Cotangent-based first association \\
13 & AQM-77 & New & Intrinsic versus relative cotangent data \\
14 & AQM-78 & Review/New & Weyl-Heisenberg kinematics from finite phase spaces \\
15 & AQM-79 & Review & Geometric field association \(\operatorname{Hom}(X,V)\) \\
16 & AQM-80 & New & Finite ring example matrix \\
17 & AQM-81 & Review/New & Residue-field site association for \(\Spec R\) \\
18 & AQM-82 & Review/New & Nilpotent-sensitive association \\
19 & AQM-83 & Review & Product rings and tensor factorization \\
20 & AQM-84 & Review/New & Infinite ring boundary cases \\
21 & AQM-85 & New & Agreement and inequivalence table \\
22 & AQM-86 & New & Degenerate and over-restrictive cases \\
23 & AQM-87 & New & Open problems and next catalogue targets \\
\bottomrule
\end{tabular}
\end{center}

\subsection{Evidence Contract}

The detailed plan fixes the first standard object families: bare finite sets,
finite examples over \(\mathbb F_3\), product rings, nilpotent Artin examples,
selected mixed-characteristic examples such as \(\mathbb Z/6\mathbb Z\), and
infinite boundary cases already present in the lab book.  Any example not yet
sourced or checked remains a question until the relevant source, convention,
derivation, or run exists.

Future comparison shards must distinguish at least four statuses: agreement
proved by an explicit map or invariant, inequivalence proved by a concrete
obstruction, degenerate collapse, and blocked/unknown due to missing
convention or source.  This shard itself proves none of those comparison
statuses.
